\section{Week 16: Final Project Presentations and Good Modeling
Practice}

\begin{quote}
\textbf{Course wrap-up} · Lecture slot: synthesis + first presentations
· Lab slot: presentations and peer feedback
\end{quote}

\begin{center}\rule{0.5\linewidth}{0.5pt}\end{center}

\subsection{Key terms}

Important terms of this lesson, with their Persian equivalents (repeated
even if they appeared in earlier weeks).

{\def\LTcaptype{none} % do not increment counter
\begin{longtable}[]{@{}
  >{\raggedright\arraybackslash}p{(\linewidth - 4\tabcolsep) * \real{0.3333}}
  >{\raggedright\arraybackslash}p{(\linewidth - 4\tabcolsep) * \real{0.3333}}
  >{\raggedright\arraybackslash}p{(\linewidth - 4\tabcolsep) * \real{0.3333}}@{}}
\toprule\noalign{}
\begin{minipage}[b]{\linewidth}\raggedright
English
\end{minipage} & \begin{minipage}[b]{\linewidth}\raggedright
\RL{\CMfa فارسی}
\end{minipage} & \begin{minipage}[b]{\linewidth}\raggedright
\RL{\CMfa توضیح کوتاه}
\end{minipage} \\
\midrule\noalign{}
\endhead
\bottomrule\noalign{}
\endlastfoot
peer feedback & \RL{\CMfa بازخورد همتا} & \RL{\CMfa نقد سازنده‌ی کار هم‌کلاسی‌ها} \\
fit / fitting & \RL{\CMfa برازش} & \RL{\CMfa تنظیم پارامترها تا مدل با داده جور
شود} \\
parameter recovery & \RL{\CMfa بازیابی پارامتر} & \RL{\CMfa آیا برازش، پارامترهای
معلوم شبیه‌سازی را برمی‌گرداند} \\
reproducibility & \RL{\CMfa بازتولیدپذیری} & \RL{\CMfa امکان به دست آوردن دوباره‌ی
همان نتایج} \\
simulation & \RL{\CMfa شبیه‌سازی} & \RL{\CMfa تولید داده‌ی مصنوعی با اجرای مدل} \\
estimate / estimator & \RL{\CMfa برآورد} / \RL{\CMfa برآوردگر} & \RL{\CMfa مقدار به‌دست‌آمده
برای پارامتر} / \RL{\CMfa روش به دست آوردن آن} \\
parameter & \RL{\CMfa پارامتر} & \RL{\CMfa کمیت قابل تنظیم مدل که بین افراد یا شرایط
فرق می‌کند} \\
likelihood & \RL{\CMfa درست‌نمایی} & \RL{\CMfa احتمال داده به‌عنوان تابعی از
پارامترها} \\
Bayes factor (BF) & \RL{\CMfa عامل بیز} (Bayes factor) & \RL{\CMfa نسبت درست‌نمایی‌های
حاشیه‌ای دو مدل} \\
bounds / constraints & \RL{\CMfa کران‌ها} / \RL{\CMfa قیدها} & \RL{\CMfa محدودیت‌های مجاز
برای مقدار پارامترها} \\
posterior & \RL{\CMfa توزیع پسین} (posterior) & \RL{\CMfa باور درباره‌ی پارامتر بعد
از دیدن داده} \\
hierarchical model & \RL{\CMfa مدل سلسله‌مراتبی (چندسطحی)} & \RL{\CMfa پارامترهای
فردی از یک توزیع گروهی می‌آیند} \\
adequacy check & \RL{\CMfa بررسی کفایت مدل} & \RL{\CMfa آیا مدل ویژگی‌های اصلی داده
را بازتولید می‌کند} \\
posterior predictive check (PPC) & \RL{\CMfa بررسی پیش‌بینی پسین} & \RL{\CMfa مقایسه‌ی
داده‌ی شبیه‌سازی‌شده از پسین با داده‌ی واقعی} \\
cross-validation (CV) & \RL{\CMfa اعتبارسنجی متقابل} (cross-validation) &
\RL{\CMfa برازش روی بخشی از داده و آزمون روی بخش کنارگذاشته} \\
AIC / BIC & \RL{\CMfa معیار اطلاعاتی آکائیکه} / \RL{\CMfa بیزی} & \RL{\CMfa برازش به‌علاوه‌ی
جریمه برای تعداد پارامتر} \\
research design & \RL{\CMfa طرح آزمایش} & \RL{\CMfa شرایط، تعداد کوشش‌ها و
نمونه‌ها} \\
model recovery & \RL{\CMfa بازیابی مدل} & \RL{\CMfa آیا مقایسه‌ی مدل، مدل مولد را
درست انتخاب می‌کند} \\
drift-diffusion model (DDM) & \RL{\CMfa مدل} drift-diffusion
\RL{\CMfa (رانش--انتشار)} & \RL{\CMfa انباشت پیوسته‌ی شواهد بین دو مرز} \\
preregistration & \RL{\CMfa پیش‌ثبت} (preregistration) & \RL{\CMfa ثبت برنامه‌ی تحلیل
پیش از دیدن داده} \\
prior & \RL{\CMfa توزیع پیشین} (prior) & \RL{\CMfa باور درباره‌ی پارامتر پیش از دیدن
داده} \\
random seed & seed \RL{\CMfa تصادفی} & \RL{\CMfa عدد شروع مولد اعداد تصادفی برای
تکرارپذیری} \\
large language model (LLM) & \RL{\CMfa مدل بزرگ زبانی} & \RL{\CMfa شبکه‌ی عصبی‌ای که
کلمه‌ی بعدی متن را پیش‌بینی می‌کند} \\
accuracy & \RL{\CMfa دقت} & \RL{\CMfa نسبت پاسخ‌های درست} \\
\end{longtable}
}

\begin{center}\rule{0.5\linewidth}{0.5pt}\end{center}

\subsection{Learning outcomes}

By the end of this week, students will be able to:

\begin{enumerate}
\def\labelenumi{\arabic{enumi}.}
\tightlist
\item
  \textbf{Communicate} a complete modeling study (question, model,
  fitting, evaluation, recovery, interpretation) clearly and within a
  time limit.
\item
  \textbf{Critique} other modeling studies constructively, using the
  course's evaluation criteria.
\item
  \textbf{Apply} a checklist of good modeling practice (reproducibility,
  transparency, robustness) to their own project report.
\item
  \textbf{Integrate} the methods of the course into one coherent
  workflow and \textbf{explain} when to use which.
\end{enumerate}

\begin{center}\rule{0.5\linewidth}{0.5pt}\end{center}

\subsection{Motivation \& big picture}

A model is only as useful as others' ability to understand, check and
build on it. Over 15 weeks, every lab has followed the same arc:
\textbf{formalize \ensuremath{\rightarrow} simulate \ensuremath{\rightarrow} fit \ensuremath{\rightarrow} evaluate \ensuremath{\rightarrow} recover \ensuremath{\rightarrow} interpret.}
This week you show that arc in your own project, and you practice the
other half of scientific modeling: reading someone else's model
critically but fairly. The session ends with the practices that make
modeling work trustworthy and reusable: reproducible code, transparent
reporting, honest statements of limitations.

\begin{center}\rule{0.5\linewidth}{0.5pt}\end{center}

\subsection{Lecture notes}

\paragraph{In plain words}

The last week has no new method. Students present their projects, and
the course is summarized as one workflow: \textbf{formalize} a theory as
a model, \textbf{simulate} it to see what it predicts, \textbf{fit} it
to data, \textbf{evaluate} it (does it fit, and does it beat
alternatives?), check that the method can \textbf{recover} the truth at
this design, and \textbf{interpret} the result in psychological terms,
including what cannot be concluded. The checklist of good practice turns
these steps into concrete things to verify in every report: reproducible
code, reported uncertainty, adequacy checks, recovery and honest
limitations.

\textbf{Suggested timing (lecture slot, 60 min):} course synthesis 15 ·
good-practice checklist 15 · first presentations 30.

\subsubsection{1. The course in one table}

{\def\LTcaptype{none} % do not increment counter
\begin{longtable}[]{@{}
  >{\raggedright\arraybackslash}p{(\linewidth - 4\tabcolsep) * \real{0.3333}}
  >{\raggedright\arraybackslash}p{(\linewidth - 4\tabcolsep) * \real{0.3333}}
  >{\raggedright\arraybackslash}p{(\linewidth - 4\tabcolsep) * \real{0.3333}}@{}}
\toprule\noalign{}
\begin{minipage}[b]{\linewidth}\raggedright
Question you face
\end{minipage} & \begin{minipage}[b]{\linewidth}\raggedright
Tools
\end{minipage} & \begin{minipage}[b]{\linewidth}\raggedright
Week
\end{minipage} \\
\midrule\noalign{}
\endhead
\bottomrule\noalign{}
\endlastfoot
How do I turn a theory into a model? & Five components; simulation;
analytic checks & 1--2 \\
How do I estimate parameters? & Least squares; maximum likelihood;
optimization (grid, Nelder--Mead, L-BFGS-B), multi-start, bounds &
3--4 \\
How certain are the estimates? & Hessian SEs, profile likelihood,
bootstrap; posterior intervals & 5, 9 \\
How do I handle many participants? & Individual fits; aggregation
artifacts; hierarchical models & 5, 13 \\
Is the model adequate? & Residuals, \(G^2\), posterior predictive checks
& 2--4, 10 \\
Which model is better? & CV, LRT, AIC/BIC, Bayes factors & 6--7, 10 \\
Can my design answer the question? & Parameter recovery, model recovery
& 7 \\
How do I get a posterior for complex models? & Grid \ensuremath{\rightarrow}
Metropolis--Hastings \ensuremath{\rightarrow} Gibbs \ensuremath{\rightarrow} NUTS (PyMC); \(\hat R\), ESS, divergences
& 9--13 \\
What do models look like in different fields? & SDT, forgetting,
practice, GCM, DDM/EZ, RL, ideal observers & 2, 5, 14--15 \\
\end{longtable}
}

\paragraph{For an engineer}

Looking back, most of the course maps onto tools you already know, and
the table below can serve as your own summary when answering students'
questions:

{\def\LTcaptype{none} % do not increment counter
\begin{longtable}[]{@{}
  >{\raggedright\arraybackslash}p{(\linewidth - 2\tabcolsep) * \real{0.5000}}
  >{\raggedright\arraybackslash}p{(\linewidth - 2\tabcolsep) * \real{0.5000}}@{}}
\toprule\noalign{}
\begin{minipage}[b]{\linewidth}\raggedright
Course topic
\end{minipage} & \begin{minipage}[b]{\linewidth}\raggedright
Engineering counterpart
\end{minipage} \\
\midrule\noalign{}
\endhead
\bottomrule\noalign{}
\endlastfoot
Random walk / DDM (Weeks 1, 14) & Sequential detection (SPRT),
first-passage times \\
SDT, ROC (Week 2) & Binary detection, \(P_D\) / \(P_{FA}\), ROC \\
Least squares, optimizers (Week 3) & LS estimation, gradient methods,
Gauss--Newton \\
Maximum likelihood, \(G^2\), LRT (Weeks 4, 7) & ML estimation,
goodness-of-fit tests, GLRT \\
Hessian SEs, bootstrap (Week 5) & Fisher information / CRLB, Monte Carlo
performance evaluation \\
CV, AIC, BIC (Weeks 6--7) & Model-order selection (AIC, MDL), validation
data \\
Parameter/model recovery (Week 7) & Testing estimators on synthetic data
with known ground truth \\
Bayesian inference, BF (Weeks 9--10) & MAP/MMSE estimation, Bayesian
detection \\
MCMC, \(\hat R\), ESS (Weeks 11--12) & Monte Carlo simulation, Markov
processes, correlated-sample statistics \\
Hierarchical models (Week 13) & Empirical Bayes, James--Stein
shrinkage \\
Delta rule, cue combination (Week 15) & LMS adaptation, MRC / Kalman
sensor fusion \\
\end{longtable}
}

What is genuinely different in cognitive modeling: the ``system'' is a
human being, the data are small and noisy, and the parameters themselves
are the scientific findings, so interpretation, uncertainty and recovery
matter more than raw performance.

\subsubsection{2. A checklist of good modeling practice}

\textbf{Before fitting} - {[} {]} The research question is stated in
terms of model parameters or model comparisons. - {[} {]} Candidate
models (including at least one simpler alternative) are specified
\emph{before} looking at the data, ideally preregistered. - {[} {]}
Priors (if Bayesian) are justified and checked with prior predictive
simulation. - {[} {]} Parameter and model recovery have been checked
\textbf{at the actual design}.

\textbf{Fitting and evaluation} - {[} {]} The likelihood matches the
data type (binomial for counts, sequential for learning, \ldots). - {[}
{]} Optimization uses multi-start and sensible bounds, and fits are
checked visually; for MCMC, \(\hat R\), ESS and divergences are
reported. - {[} {]} Absolute fit (residuals, \(G^2\) or PPC) \emph{and}
relative fit (CV/IC/BF) are reported. - {[} {]} Uncertainty is reported
for every interpreted parameter.

\textbf{Reporting and reproducibility} - {[} {]} All code and (where
allowed) data are shared; random seeds are fixed. - {[} {]} Software
versions are recorded (e.g.,
\nolinkurl{conda env export > environment.yml}). - {[} {]} One command
or notebook regenerates every number and figure in the report. - {[} {]}
Limitations are stated: what the design cannot distinguish, which
parameters are not recoverable, which assumptions were not tested.

\textbf{Recording the computing environment} (include in every project
notebook):

\begin{quote}
\textbf{Code · \RL{\CMfa کد:}} \nolinkurl{code/week16.py}, cell 1 · \RL{\CMfa سلول}
1
\end{quote}

\subsubsection{3. Presentation format}

\begin{itemize}
\tightlist
\item
  \textbf{Length:} 8 minutes (individual) or 10 minutes (pairs), plus 4
  minutes of questions. Strict timekeeping.
\item
  \textbf{Suggested structure} (one slide each, roughly):

  \begin{enumerate}
  \def\labelenumi{\arabic{enumi}.}
  \tightlist
  \item
    The question and why it matters.
  \item
    The model(s): equations and the psychological meaning of every
    parameter.
  \item
    Data and design.
  \item
    Fitting method and a fit figure (data + predictions).
  \item
    Evaluation: model comparison \emph{and} an adequacy check.
  \item
    Recovery: can the design answer the question?
  \item
    Interpretation and limitations.
  \end{enumerate}
\item
  \textbf{Pairs:} both members present, and the Bayesian/MCMC or
  hierarchical component must be shown.
\item
  \textbf{Schedule:} posted one week in advance; the order is drawn at
  random. Sessions run across the lecture and lab slots (about 15--18
  presentations in 3 hours, depending on enrolment).
\end{itemize}

\subsubsection{4. Giving and receiving feedback}

Each student completes a \textbf{peer-feedback form} for three assigned
presentations (below). Feedback is anonymous to the presenter, visible
to the instructor, and counts toward participation (lab check-off). Good
feedback is \textbf{specific} (``the recovery plot shows \(\beta\) is
not identifiable, so the group difference in \(\beta\) should not be
interpreted'') rather than general (``nice talk''), and
\textbf{constructive} (it suggests a concrete improvement).

\textbf{Peer-feedback form}

{\def\LTcaptype{none} % do not increment counter
\begin{longtable}[]{@{}
  >{\raggedright\arraybackslash}p{(\linewidth - 4\tabcolsep) * \real{0.3333}}
  >{\raggedright\arraybackslash}p{(\linewidth - 4\tabcolsep) * \real{0.3333}}
  >{\raggedright\arraybackslash}p{(\linewidth - 4\tabcolsep) * \real{0.3333}}@{}}
\toprule\noalign{}
\begin{minipage}[b]{\linewidth}\raggedright
Item
\end{minipage} & \begin{minipage}[b]{\linewidth}\raggedright
Rating (1--4)
\end{minipage} & \begin{minipage}[b]{\linewidth}\raggedright
Comment
\end{minipage} \\
\midrule\noalign{}
\endhead
\bottomrule\noalign{}
\endlastfoot
The question was clear and the model matched it & & \\
The model's parameters were explained in psychological terms & & \\
The fit was shown and evaluated (adequacy + comparison) & & \\
Recovery or identifiability was addressed & & \\
Conclusions followed from the results, and limitations were stated &
& \\
\textbf{One thing that was particularly well done} & & \\
\textbf{One concrete suggestion for the written report} & & \\
\end{longtable}
}

\paragraph{For the instructor}

\begin{itemize}
\tightlist
\item
  \textbf{Grade with the rubric, comment with the checklist.} During
  presentations, keep the good-practice checklist next to you and note
  one strength and one concrete improvement per team; this makes the
  written feedback fast and fair.
\item
  \textbf{Questions to ask presenters} (they map onto the course): Which
  parameter answers your question? How do you know it is recoverable?
  Does the model fit in absolute terms, not only better than the
  alternative? What can your design \emph{not} tell us?
\item
  \textbf{Use the worked critique} (fictional Q-learning study) as a
  model for how students should give feedback: specific, tied to a week
  of the course, and constructive.
\item
  \textbf{Close the course with the one-table summary.} Ask students
  which row they found hardest, and point them to the corresponding week
  for exam revision.
\end{itemize}

\begin{center}\rule{0.5\linewidth}{0.5pt}\end{center}

\subsection{Worked example: critiquing a modeling study}

\textbf{Summary of a (fictional) study.} \emph{``We fitted a Q-learning
model to 30 patients and 30 controls who performed 60 trials of a
two-armed bandit task with stable reward probabilities. Each participant
was fitted by MLE. Patients had significantly lower learning rates
(\(t(58) = 2.4\), \(p = .02\)). We conclude that the disorder impairs
reward learning.''}

\textbf{Critique using the checklist.}

\begin{enumerate}
\def\labelenumi{\arabic{enumi}.}
\tightlist
\item
  \textbf{Recovery at this design (Week 7, Week 15 lab).} With 60 trials
  and \emph{stable} reward probabilities, the learning rate is weakly
  identifiable, and \(\alpha\) and \(\beta\) trade off. Without a
  recovery study, a difference in \(\hat\alpha\) could reflect a
  difference in choice noise (\(\beta\)) or just estimation noise.
  \emph{Suggestion:} report parameter recovery at this design, and
  include reversals.
\item
  \textbf{Model comparison (Week 15, Part D).} Did Q-learning beat
  simpler strategies (WSLS, perseveration) in both groups? If patients
  are better described by a different strategy, comparing their
  \(\alpha\) with controls' is not meaningful. \emph{Suggestion:} report
  AIC/BIC or CV per group, with model recovery.
\item
  \textbf{Second-stage t-test on point estimates (Weeks 5, 13).} It
  ignores the uncertainty of individual estimates, which is large with
  60 trials. \emph{Suggestion:} a hierarchical model with a group effect
  on \(\alpha\), reporting its posterior.
\item
  \textbf{Adequacy (Weeks 4, 10).} No posterior predictive check or
  learning-curve plot is shown. Does the model reproduce each group's
  learning curve?
\item
  \textbf{Interpretation.} ``Impairs reward learning'' is a strong
  process claim. A lower \(\hat\alpha\) is also compatible with more
  exploration, lower motivation (lower \(\beta\)), or different priors
  on initial values. \emph{Suggestion:} temper the conclusion, or add a
  design that separates these accounts.
\end{enumerate}

\textbf{Overall:} a reasonable first study whose conclusion is not yet
supported. Three additions (recovery, model comparison, hierarchical
estimation) would make it convincing. Notice that each point maps
directly onto one week of this course.

\begin{center}\rule{0.5\linewidth}{0.5pt}\end{center}

\subsection{Lab (2 hours): presentations and peer feedback}

\textbf{Goals:} present your project; give written feedback on three
presentations; note the feedback you receive for your final report.

\textbf{Logistics:} - Upload slides (PDF) to the course page
\textbf{before} the session starts. - Bring your laptop with your
project notebook in case of questions about the code. - Fill in the
peer-feedback forms during the session (paper or online form) and submit
them by the end of the day.

\textbf{After the presentations:} the instructor gives brief oral
feedback to each team. Written feedback (instructor + peers) is returned
within 3 days so it can be used in the final report.

\begin{center}\rule{0.5\linewidth}{0.5pt}\end{center}

\subsection{Connection to AI \& LLMs}

\textbf{Good practice when AI is part of your modeling work.} AI tools
now enter cognitive modeling in three roles, and each needs the course's
habits of checking: 1. \textbf{As a coding assistant.} Generated code
can look right and be wrong. Test it the modeling way: simulate data
with known parameters and check that the code recovers them (parameter
recovery as a unit test). 2. \textbf{As a ``simulated participant''.}
Some studies give LLMs psychology experiments and compare the answers
with human data, or even propose LLM-based models of human behavior
(e.g., Binz \& Schulz, 2023; Dillion et al., 2023 {[}verify{]}).
Evaluate such claims like any model: generalization to new tasks,
comparison with simpler cognitive models, mimicry, and interpretability
(Week 15). 3. \textbf{As an object of study.} If you test an LLM, record
the model name and version, the date, the exact prompts, the temperature
and the number of samples. Without these, the study cannot be
reproduced.

The code below tests a pre-made SDT function ``written by an
assistant''. It looks plausible, and it even passes the case \(c = 0\),
but a recovery test shows that it gets the sign of \(c\) wrong.

\begin{quote}
\textbf{Code · \RL{\CMfa کد:}} \nolinkurl{code/week16.py}, cell 2 · \RL{\CMfa سلول}
2
\end{quote}

\textbf{Expected output:}

\Needspace{6\baselineskip}
\begin{Verbatim}
true (d', c) = (1.5, 0.3) -> recovered (1.50, -0.30)  FAIL
true (d', c) = (1.0, -0.4) -> recovered (1.00, 0.40)  FAIL
true (d', c) = (2.0, 0.0) -> recovered (2.00, 0.00)  OK
\end{Verbatim}

A liberal observer would be reported as conservative, which would
reverse the conclusion of a bias study. The test with \(c = 0\) alone
would not have caught it.

\textbf{Try it (no API needed).} - \ensuremath{\star} Fix the function and rerun the
tests. - \ensuremath{\star}\ensuremath{\star} Write a similar recovery test for the EZ-diffusion function
of Week 14, using the simulator from the same lab. - \ensuremath{\star}\ensuremath{\star} For your own
project: add one paragraph to the reproducibility statement that
describes how AI tools were used and how their output was verified (see
\nolinkurl{assessments/project_brief.md}, section 6).

\emph{Optional reading:} Binz, M., \& Schulz, E. (2023). Using cognitive
psychology to understand GPT-3. \emph{PNAS, 120}(6), e2218523120.
Dillion, D., Tandon, N., Gu, Y., \& Gray, K. (2023). Can AI language
models replace human participants? \emph{Trends in Cognitive Sciences,
27}, 597--600 {[}verify{]}.

\begin{center}\rule{0.5\linewidth}{0.5pt}\end{center}

\subsection{Common mistakes \& misconceptions (in project presentations
and reports)}

\begin{enumerate}
\def\labelenumi{\arabic{enumi}.}
\tightlist
\item
  \textbf{Too much on the method, too little on the meaning.} Every
  parameter and every result needs a psychological interpretation.
\item
  \textbf{Showing a fit statistic without a fit figure.} Always plot
  data against predictions.
\item
  \textbf{Model comparison without an adequacy check,} or vice versa.
\item
  \textbf{Interpreting parameters that were not shown to be
  recoverable.}
\item
  \textbf{Overclaiming:} ``the model is correct'', ``the mind is
  Bayesian''. State what the evidence supports and what it does not.
\item
  \textbf{Irreproducible results:} missing seeds, unrecorded package
  versions, notebook cells run out of order.
\item
  \textbf{Running over time.} Practice the talk; one idea per slide.
\end{enumerate}

\begin{center}\rule{0.5\linewidth}{0.5pt}\end{center}

\subsection{Check your understanding (integrative review questions for
the final exam)}

\textbf{1.} A colleague has accuracy and RT data from a speed/accuracy
manipulation and wants to know whether the manipulation affects caution
or evidence quality. Outline a full modeling plan, naming the model, the
fitting method, the evaluation, and the recovery check. \textbf{Answer}
Model: drift-diffusion (EZ-diffusion for a first pass, or the full DDM
by MLE / hierarchical Bayes). Fit per participant and condition. Compare
models in which the manipulation affects only \(a\), only \(v\), or both
(AIC/BIC or a hierarchical model with condition effects; PPC on RT
quantiles). Recovery: simulate data at the actual design with each
model, and check that parameters and the ``which parameter changed''
conclusion are recovered. Interpret the condition effects on \(a\) and
\(v\) with uncertainty.

\textbf{2.} When would you choose a Bayes factor over cross-validation
for comparing two models, and vice versa? \textbf{Answer} Bayes factor:
when the question is which model is more probable given the data
(evidence, including support for a null), when meaningful priors are
available, and with a prior sensitivity analysis. CV: when the goal is
predictive accuracy for new data of a specific kind, priors are hard to
justify, or the BF is too sensitive to the priors.

\textbf{3.} Explain how a parameter can be estimated very precisely
(narrow CI or posterior) and yet be meaningless. \textbf{Answer}
Precision only reflects the curvature of the likelihood under the
assumed model. If the model is misspecified (fails PPC/\(G^2\)), if the
parameter trades off with another so that its meaning depends on an
arbitrary constraint, or if the optimizer or sampler failed (local
minimum, bound, non-convergence), the precise estimate does not measure
the intended construct.

\textbf{4.} Why do hierarchical models often give better individual
estimates than separate individual fits? \textbf{Answer} They share
information across participants through the group distribution,
shrinking noisy individual estimates toward the group mean in proportion
to their uncertainty. This reduces estimation variance more than it adds
bias, so estimates are on average closer to the truth.

\textbf{5.} Name three different ways in which this course used
simulation, and why each matters. \textbf{Answer} (i) Exploring a
model's predictions before fitting (what can and can't it produce). (ii)
Parameter and model recovery (can the design and method answer the
question). (iii) Bootstrap and predictive checks (uncertainty and
adequacy). Further: Monte Carlo and MCMC for computing posteriors.

\textbf{6.} A model has the best AIC among four candidates but a
posterior predictive check shows a systematic misfit at long delays.
What do you report and conclude? \textbf{Answer} Report both. The model
is the best of the candidates for prediction, but it is not adequate in
absolute terms: none of the candidates captures the long-delay behavior.
Conclusions should be restricted accordingly, and a model revision
targeting that feature is the next step.

\begin{center}\rule{0.5\linewidth}{0.5pt}\end{center}

\subsection{Assignment}

\subsubsection{Final project: written report and code}

\begin{itemize}
\tightlist
\item
  \textbf{Due:} one week after this session, and in any case before the
  final exam (exact date on the course page).
\item
  \textbf{Submission:} report (PDF, max. 10 pages for individuals, 14
  for pairs, excluding references and appendix) plus a code archive or
  repository link that reproduces all results.
\item
  \textbf{Grading:} see the project brief and rubric
  (\nolinkurl{assessments/project_brief.md}). The report and code are
  worth 3 of the 6 project points; the proposal (1) and presentation (2)
  are already graded.
\item
  \textbf{Use the feedback} from today: the rubric rewards visible
  responses to presentation feedback.
\end{itemize}

\subsubsection{Final exam}

Written, cumulative (emphasis on Chapters 4--6), formula sheet provided,
no laptops; worth 3 of 20 points. See
\nolinkurl{assessments/final_exam.md} for the format and a sample paper.

\begin{center}\rule{0.5\linewidth}{0.5pt}\end{center}

\subsection{Further reading}

\begin{itemize}
\tightlist
\item
  \textbf{Farrell \& Lewandowsky (2018)}: the chapter on using models in
  psychology (the role of models in theory development and common
  pitfalls) {[}verify chapter number{]}.
\item
  \textbf{Lee \& Wagenmakers (2014)}: browse the case studies in Part IV
  for ideas on how to extend your project.
\item
  \textbf{Martin, Kumar \& Lao (2021), chapter 9.} End-to-end Bayesian
  workflows, a good model for structuring a report.
\item
  \emph{Optional, outside the course list:} Wilson, R. C., \& Collins,
  A. G. E. (2019). Ten simple rules for the computational modeling of
  behavioral data. \emph{eLife, 8}, e49547.
\end{itemize}

\textbf{Thank you for a great course, and keep modeling!}
